%\documentclass[10pt,letterpaper]{preprint}
\documentclass[10pt,letterpaper]{article}

\usepackage{psepress}
\addbibresource{refs.bib}

% Change `preprint` above back to `article` to build the final two-column version.
% Edit this metadata block when starting a new paper.
\renewcommand{\PaperTitle}{Place Your Paper Title Here with Title Capitalization that May Run Multiple Lines}
\renewcommand{\PaperAuthors}{%
Firstname Q.~Lastname\textsuperscript{a},
Givenname R.~S.~Surname Jr.\textsuperscript{ab},
and Personwith MULTIWORD LASTNAME\textsuperscript{b*}}
\renewcommand{\PaperAffiliations}{%
\textsuperscript{a} Provincial University, Department of Chemical Engineering, City, Province, Country\\
\textsuperscript{b} State College, Engineering Department, City, State, Country\\
\textsuperscript{*} Corresponding Author: \href{mailto:author.email@someuniversity.ca}{author.email@someuniversity.ca}}
\renewcommand{\PaperAbstract}{%
Type your abstract text here. Please limit to 250 words. This block is already formatted with the
correct abstract styling for the template.
Lorem ipsum dolor sit amet, consectetur adipiscing elit, sed do eiusmod tempor incididunt ut labore
et dolore magna aliqua. Ut enim ad minim veniam, quis nostrud exercitation ullamco laboris nisi ut
aliquip ex ea commodo consequat. Duis aute irure dolor in reprehenderit in voluptate velit esse
cillum dolore eu fugiat nulla pariatur. Excepteur sint occaecat cupidatat non proident, sunt in culpa
qui officia deserunt mollit anim id est laborum.}
\renewcommand{\PaperKeywords}{%
Choose an item., Choose an item., Choose an item., Choose an item., Choose an item. additional
keywords separated by commas}
\renewcommand{\PaperCopyrightText}{%
\textcopyright\ 2027 by the authors. Licensed to PSEcommunity.org and PSE Press. This is an open
access article under the creative commons CC-BY-SA licensing terms. Credit must be given to creator
and adaptations must be shared under the same terms. See
\url{https://creativecommons.org/licenses/by-sa/4.0/}.}

\renewcommand{\HeaderArticleType}{Original Research Article}
\renewcommand{\HeaderReviewType}{Peer Reviewed Conference Proceeding}
% Conference selector. Built-in choices are:
%   escape-37-2027, pse-2027, and focapo-cpc-2027.
\PSESelectConference{escape-37-2027}
% For another conference, replace the selector above with, for example:
% \PSESetConference{Conference Name}{City, Country, 6-9 June 2027}
\renewcommand{\FooterJournalInfo}{Syst Control Trans V:XXXX-YYYY (2027)}
\renewcommand{\FirstPageFooterID}{[DOI\_DoNotChange]}
\renewcommand{\BodyPageFooterID}{[LAPSE\_DoNotChange]}

\begin{document}
\psemaketitle

\section{this is header 1}
\subsection{This is Header 2}
Regular text is already formatted for the body style used by this template. Do not add blank lines
between paragraphs for visual spacing. Paragraphs indent automatically and will hyphenate
automatically.

\subsubsection{This is header 3}
Use the provided sectioning commands for headings. The spacing before and after headings is already
set, so do not add extra blank paragraph space manually.

\paragraph{This is header 4}
Some text that follows header 4.

\section{Figures and Tables in the text}
\subsection{One Column-Width Figures}
The standard column width is 8.5 cm in the final article layout. In review mode, \texttt{\string\columnwidth}
expands to the full single-column text width automatically. For a normal figure, use a
\texttt{figure} environment and keep the graphic at or below \texttt{\string\columnwidth}. The
example below uses \texttt{[htp!]} placement so the sample stays near the surrounding explanatory
text while still allowing normal float behavior.

\begin{figure}[htp!]
  \centering
  \includegraphics[width=\columnwidth]{one-column-figure.png}
  \caption{Example 1-column figure with a short caption. \textit{Max width} is 8.5 cm in the final article layout.}
\end{figure}

\subsection{One-Column Tables}
Use a \texttt{table} environment for one-column tables. The sample below also uses \texttt{[htp!]}
so it stays near the nearby text without forcing exact placement.

\begin{table}[htp!]
  \caption{Caption for the table. Multi-line table captions will be fully justified.}
  \centering
  \renewcommand{\arraystretch}{1.08}
  \begin{tabularx}{\columnwidth}{@{} l X r @{}}
    \toprule
    \textbf{Stream} & \textbf{Description} & \textbf{Flow} \\
    \midrule
    Feed & MeOH/EtOH/BuOH blend & 11.25 kmol/hr \\
    Side Draw & Ethanol-rich cut & 115.00 kmol/hr \\
    Distillate & Methanol-rich overhead & 21.50 kmol/hr \\
    Bottoms & Butanol-rich bottoms & 110.00 kmol/hr \\
    \bottomrule
  \end{tabularx}
\end{table}

\subsection{Two Column (Page-Width) Tables and Figures}
The printable page width is 17.5 cm, but images and tables should be limited to 17 cm wide in the
final article layout. In LaTeX, use \texttt{figure*} or \texttt{table*} only when the content truly
needs both columns. In the review build there is only one column, so \texttt{figure} and
\texttt{table} should be preferred.

\section{How to CITE REFERENCES IN THE TEXT}
Text citations should use numbers in square brackets [1] and be placed at the end of the phrase that
is most appropriate in the context of the sentence, which may or may not be at the end of the
sentence. In this template, bibliography data lives in \texttt{refs.bib} and citations use
\texttt{\string\cite\{key\}}. You can cite ranges of papers [1-5] or sets of papers [3,5,8].

\section{Other Formatting instructions}
\subsection{Equations in the Text}
Equations should look like the following. Use standard displayed equations and let LaTeX align the
equation numbers flush right. Do not add additional space above or below the equation; the
environment handles that automatically.

\begin{equation}
  f_x = \mu \int_0^\infty \frac{1}{x}\,dx
\end{equation}

\begin{equation}
  f_x = \alpha\beta \int_0^\infty \frac{1}{x}\,dx + y^2
\end{equation}

You can use LaTeX math directly here. Here is an example of a matrix:

\begin{equation}
  A =
  \left\{
  \begin{matrix}
    a & b & c \\
    d & e & f \\
    g & h & j
  \end{matrix}
  \right\}
\end{equation}

\subsection{Code}
Use \texttt{\string\texttt\{\}} for short inline code such as a function name. For code blocks, use
a monospace environment as demonstrated below.

\begin{lstlisting}[style=psecode]
for i in range(NPoints):
    T.append(600 + (400) * (i/(NPoints-1)))
    P.append(  2 +  (50) * (i/(NPoints-1)))

random.shuffle(P)
\end{lstlisting}

\subsection{Language and Spelling}
We prefer US English. Other forms are acceptable.

\subsection{Headers, Footers, and Copyright}
Do not change the package-provided page geometry, header/footer settings, or the closing copyright
block in the final article layout. The review build removes the branded conference header and uses a
plain one-column preprint layout automatically.

\subsection{Definitions, Theorems, and Block Quotes}
Use the \texttt{psedefinition} environment included in this template for theorem-like statements,
definitions, and block quotations. Give it a short title such as ``Theorem 1'' or ``Definition'' as
its argument.

\begin{psedefinition}{Theorem 1}
Then this text is where your theorem is declared, and you can describe it across multiple lines.
\end{psedefinition}

\begin{psedefinition}{Term Being Defined}
This works well for definitions too.
\end{psedefinition}

\section{DIGITAL supplementary material}
You are encouraged to submit digital supplementary materials to the Living Archive for Process
Systems Engineering at \url{http://PSEcommunity.org/LAPSE}. You can then refer to your digital
materials by using the LAPSE ID in this section. Delete this section if not used.

\section{ACKNOWLEDGEMENTS}
Please acknowledge your funders in this section, with grant or funding numbers as appropriate. Note
that most funding agencies require this.

\section{DELCARATION OF USE OF AI}
Describe how generative AI was used in preparing the manuscript. If AI was used in the research
methods, describe that in the main paper instead (for example, in the methods section). AI is not an
author.

\section{author identifiers}
Author ORCIDs:

Lastname FQ: 0000-0001-2345-6789

Surname GRS Jr: 0000-0001-2345-6789

Multiword Lastname P: 0000-0001-2345-6789

% Remove \nocite{*} once you want the bibliography to include only cited entries.
\nocite{*}
\pseprintreferences
\pseprintcopyright

\end{document}
